\documentclass{article}
\usepackage[T1]{fontenc}
\usepackage{lmodern}
\usepackage{graphicx}
\usepackage{amsmath,amssymb}
\usepackage[a4paper,margin=2cm]{geometry}
\usepackage[absolute,overlay]{textpos}
\setlength{\TPHorizModule}{1mm}
\setlength{\TPVertModule}{1mm}
\usepackage[indonesian]{babel}
\usepackage{hyperref}
\setlength{\emergencystretch}{2em}
% unit: Halaman 87

\begin{document}

% === Halaman 87 (llm) ===

\section*{Kelebihan Salsa20}

\begin{enumerate}
    \item \textbf{Sangat cepat}\
    Salsa20 hanya menggunakan operasi yang sangat sederhana (AXR):
    \begin{itemize}
        \item penjumlahan (\textit{add}) modulo $2^{32}$,
        \item XOR,
        \item rotasi bit.
    \end{itemize}

    \item \textbf{Memiliki keamanan yang baik}\
    Salsa20 dirancang dengan 20 putaran untuk memberikan difusi yang kuat. Maksudnya, perubahan kecil pada \textit{key, nonce, counter}, akan menghasilkan perubahan yang besar pada keystream.
\end{enumerate}

\end{document}
